
\documentclass[11pt]{article}

\usepackage[margin=1in]{geometry}
\usepackage[T1]{fontenc}
\usepackage[utf8]{inputenc}
\usepackage{lmodern}
\usepackage{microtype}
\usepackage{amsmath,amssymb,amsthm,mathtools}
\usepackage{enumitem}
\usepackage{booktabs,longtable,array}
\usepackage{xcolor}
\usepackage{hyperref}
\usepackage{titlesec}

\hypersetup{
  colorlinks=true,
  linkcolor=blue!50!black,
  citecolor=blue!50!black,
  urlcolor=blue!50!black,
  pdftitle={Notes on Erd\H{o}s Problem \#885},
  pdfauthor={OpenAI / user notes synthesis}
}

\titleformat{\section}{\Large\bfseries}{\thesection.}{0.6em}{}
\titleformat{\subsection}{\large\bfseries}{\thesubsection.}{0.6em}{}
\titleformat{\subsubsection}{\normalsize\bfseries}{\thesubsubsection.}{0.6em}{}

\newtheoremstyle{plainbreak}
  {\topsep}{\topsep}
  {\itshape}{}
  {\bfseries}{.}
  { }{\thmname{#1}\thmnumber{ #2}\thmnote{ (#3)}\newline}
\theoremstyle{plainbreak}
\newtheorem{theorem}{Theorem}[section]
\newtheorem{proposition}[theorem]{Proposition}
\newtheorem{lemma}[theorem]{Lemma}
\newtheorem{corollary}[theorem]{Corollary}

\newtheoremstyle{defbreak}
  {\topsep}{\topsep}
  {}{}
  {\bfseries}{.}
  { }{\thmname{#1}\thmnumber{ #2}\thmnote{ (#3)}\newline}
\theoremstyle{defbreak}
\newtheorem{definition}[theorem]{Definition}
\newtheorem{computation}[theorem]{Verified Computation}
\newtheorem{observation}[theorem]{Observation}
\newtheorem{remark}[theorem]{Remark}
\newtheorem{example}[theorem]{Example}
\newtheorem{question}[theorem]{Open Question}

\DeclareMathOperator{\Jac}{Jac}
\DeclareMathOperator{\Aut}{Aut}
\DeclareMathOperator{\CR}{CR}
\DeclareMathOperator{\ord}{ord}
\DeclareMathOperator{\rank}{rank}

\newcommand{\Dset}{D}
\newcommand{\Fset}{F}
\newcommand{\Sig}{\Sigma}
\newcommand{\Aset}{A}
\newcommand{\Bset}{B}
\newcommand{\Mset}{M}
\newcommand{\Q}{\mathbf{Q}}
\newcommand{\Z}{\mathbf{Z}}
\newcommand{\PP}{\mathbf{P}}

\begin{document}

\begin{titlepage}
\centering
{\Large \textbf{Notes on Erd\H{o}s Problem \#885}}\\[1.25em]
{\large Exact Reductions, Verified Computations, Packet Laws, Row Complexes, and Seed-Centered Deformations}\\[2em]
\begin{minipage}{0.88\textwidth}
\small
These notes consolidate the exact statements, proofs, verified computations, and structural reductions that emerged across the full working session.  The aim is not to present a final proof of the full conjecture, but to collect every mathematically valuable theorem, proposition, computation, and obstruction in one coherent \LaTeX{} source.  Statements are deliberately separated into three classes:
\begin{enumerate}[label=\textbf{(\arabic*)},itemsep=0.4em]
    \item \textbf{Theorems / propositions / lemmas}: exact derivations proved in the notes.
    \item \textbf{Verified computations}: exact finite computations, exact enumerations, or symbolic calculations recorded in the notes and rechecked during synthesis.
    \item \textbf{Open directions / observations}: structurally important but unresolved consequences.
\end{enumerate}
The document is organized so that each route --- finite fibers, anchored packets, divisor grids, row complexes, square-translate geometry, and the non-split seed-centered deformation --- can be read independently.
\end{minipage}
\vfill
{\large Prepared from user research notes and session derivations}\\[0.5em]
{\large \today}
\end{titlepage}

\tableofcontents
\bigskip

\section{Problem statement and basic notation}

\begin{definition}
For a positive integer $n$, define its \emph{factor-difference set}
\[
\Dset(n):=\{\,|a-b| : n=ab,\ a,b\in\Z_{>0}\,\}.
\]
\end{definition}

The problem studied throughout is the following.

\begin{quote}
\textbf{Erd\H{o}s Problem \#885.}
For every $k\ge 1$, do there exist positive integers
\[
N_1<\cdots<N_k
\]
and positive integers
\[
d_1<\cdots<d_k
\]
such that
\[
d_i\in \Dset(N_j)\qquad (1\le i,j\le k)?
\]
Equivalently, can one find $k$ integers with at least $k$ common factor differences?
\end{quote}

When a finite set $S\subset \Z_{>0}$ is fixed, we write
\[
\Fset(S):=\{n\ge 1:S\subseteq \Dset(n)\}.
\]
Likewise, when a finite set $T\subset \Z_{>0}$ is viewed as a seed of common differences, we write
\[
\Sigma(T):=\{n\ge 1:T\subseteq \Dset(n)\}.
\]

\begin{proposition}[square-discriminant reformulation]
For positive integers $d,n$, one has
\[
d\in \Dset(n)\iff d^2+4n \text{ is a square}.
\]
More precisely,
\[
d\in \Dset(n)\iff \exists\,m\in\Z_{>0}\text{ such that }m^2-d^2=4n.
\]
\end{proposition}

\begin{proof}
If $d\in \Dset(n)$, then $n=a(a+d)$ for some $a\in\Z_{>0}$.  Set $m:=2a+d$. Then
\[
m^2-d^2=(2a+d)^2-d^2=4a(a+d)=4n.
\]
Conversely, if $m^2-d^2=4n$, then
\[
n=\frac{(m-d)(m+d)}{4}.
\]
Since $m^2\equiv d^2\pmod 4$, the integers $m-d$ and $m+d$ have the same parity, hence
\[
a:=\frac{m-d}{2},\qquad b:=\frac{m+d}{2}
\]
are positive integers with $ab=n$ and $b-a=d$.
\end{proof}

\begin{remark}
In the square-translate language, the whole problem asks whether for every $k$ one can find sets
\[
X=\{4N_1,\dots,4N_k\},\qquad Y=\{d_1^2,\dots,d_k^2\}
\]
with $|X|=|Y|=k$ and
\[
X+Y\subseteq \{\text{squares}\}.
\]
This viewpoint reappears later in the quartic, elliptic, and hyperelliptic reformulations.
\end{remark}

\section{Fixed finite fibers and pairwise algebra}

\subsection{Exact parametrization of $\Fset(\{d,e\})$}

\begin{theorem}[exact parametrization of a pair-fiber]
Let $0<d<e$, and put $\Delta:=e^2-d^2$. Then $n\in \Fset(\{d,e\})$ if and only if there exist integers $r,s$ such that
\[
rs=\Delta,\qquad r<s,\qquad r\equiv s\pmod 2,
\]
and, writing
\[
x:=\frac{s-r}{2},\qquad y:=\frac{s+r}{2},
\]
one has
\[
x>d,\qquad x\equiv d\pmod 2,\qquad n=\frac{x^2-d^2}{4}.
\]
Equivalently,
\[
\Fset(\{d,e\})\longleftrightarrow
\left\{(r,s):
\begin{array}{l}
rs=e^2-d^2,\\
r<s,\ r\equiv s\pmod 2,\\
\frac{s-r}{2}>d,\ \frac{s-r}{2}\equiv d\pmod 2
\end{array}\right\}.
\]
\end{theorem}

\begin{proof}
By Proposition 1.2, $n\in \Fset(\{d,e\})$ if and only if there exist integers $x<y$ such that
\[
x^2=d^2+4n,\qquad y^2=e^2+4n.
\]
Subtracting gives
\[
(y-x)(y+x)=e^2-d^2=\Delta.
\]
Setting
\[
r:=y-x,\qquad s:=y+x
\]
gives the required factorization with $r\equiv s\pmod 2$. Conversely, from such a factorization one recovers $x,y$, hence $n=(x^2-d^2)/4$.
\end{proof}

\begin{corollary}[finite-fiber lemma]
If $S\subset \Z_{>0}$ has at least two distinct elements, then $\Fset(S)$ is finite.
\end{corollary}

\begin{proof}
Choose distinct $d,e\in S$. Then
\[
\Fset(S)\subseteq \Fset(\{d,e\}),
\]
and Theorem 2.1 shows the latter is finite.
\end{proof}

\begin{remark}
This single lemma already rules out any construction with a \emph{fixed} common-difference set $S$ and infinitely many supporting integers. Any general construction for all $k$ must let the common differences vary with the tuple.
\end{remark}

\subsection{A sharp count formula in the same-parity case}

\begin{proposition}[same-parity count formula]
Assume $d$ and $e$ have the same parity. Write
\[
d=B-A,\qquad e=B+A,\qquad 0<A<B.
\]
Then $e^2-d^2=4AB$, and
\[
|\Fset(\{d,e\})|=
\#\left\{u\mid AB:\ u<A,\ \frac{AB}{u}-u\equiv d\pmod 2\right\}.
\]
\end{proposition}

\begin{proof}
In Theorem 2.1 write $r=2u$ and $s=2v$. Then $uv=AB$ and
\[
x=\frac{s-r}{2}=v-u.
\]
The condition $x>d=B-A$ is equivalent to $u<A$: indeed, if $u<A$, then $v=AB/u>B$, hence $v-u>B-A$; and if $u\ge A$, then $v\le B$, hence $v-u\le B-A$. The parity condition becomes $v-u\equiv d\pmod 2$, giving the formula.
\end{proof}

\begin{corollary}[odd-divisor largeness criterion]
Under the same hypotheses, each odd divisor $q>1$ of $A$ produces an element of $\Fset(\{d,e\})$ by taking
\[
u=\frac{A}{q},\qquad v=Bq,\qquad
n_q=\frac{(Bq-A/q)^2-d^2}{4}.
\]
Consequently,
\[
|\Fset(\{d,e\})|\ge \tau_{\mathrm{odd}}(A)-1.
\]
In particular, if $A$ is odd, then
\[
|\Fset(\{d,e\})|\ge \tau(A)-1.
\]
\end{corollary}

\begin{proof}
If $q\mid A$ is odd, then $u=A/q$ and $v=Bq$ satisfy $uv=AB$ and $u<A$. Because $q$ is odd,
\[
v-u=Bq-\frac{A}{q}\equiv B-A=d\pmod 2,
\]
so Proposition 2.4 applies.
\end{proof}

\section{Verified small examples and exact minima}

\begin{computation}[the first nontrivial pair]
The lexicographically minimal pair $(N_1,N_2)$ with two common positive factor differences is
\[
(12,42),
\]
with
\[
\Dset(12)\cap \Dset(42)=\{1,11\}.
\]
Moreover,
\[
\Fset(\{1,11\})=\{12,42,210\}.
\]
\end{computation}

\begin{proof}[Verification]
Indeed
\[
12=3\cdot 4=1\cdot 12,\qquad 42=6\cdot 7=3\cdot 14,
\]
so both $1$ and $11$ occur in each factor-difference set.  The factorization
\[
11^2-1^2=120
\]
has exactly the relevant same-parity decompositions yielding $n=12,42,210$.
\end{proof}

\begin{computation}[the first nontrivial triple]
The lexicographically minimal triple $(N_1,N_2,N_3)$ with three common positive factor differences is
\[
(112,952,3240),
\]
with
\[
\Dset(112)\cap \Dset(952)\cap \Dset(3240)=\{6,54,111\}.
\]
Moreover,
\[
\Fset(\{6,54,111\})=\{112,952,3240\}.
\]
\end{computation}

\begin{proof}[Verification]
One checks
\[
112=8\cdot 14=2\cdot 56=1\cdot 112,
\]
\[
952=28\cdot 34=14\cdot 68=8\cdot 119,
\]
\[
3240=54\cdot 60=18\cdot 72=24\cdot 135,
\]
which yield differences $6,54,111$ in each case.  The exact fiber statement follows from finite enumeration of $\Fset(\{6,54,111\})$ via Theorem 2.1.
\end{proof}

\begin{computation}[a non-endpoint triple]
There is a triple
\[
(120,528,4488)
\]
with
\[
\Dset(120)\cap \Dset(528)\cap \Dset(4488)=\{2,37,58\}.
\]
Equivalently,
\[
\Fset(\{2,37,58\})=\{120,528,4488\}.
\]
\end{computation}

\begin{proof}[Verification]
Indeed
\[
120=10\cdot 12=3\cdot 40=2\cdot 60,
\]
\[
528=22\cdot 24=11\cdot 48=8\cdot 66,
\]
\[
4488=66\cdot 68=51\cdot 88=44\cdot 102.
\]
The common differences are exactly $2,37,58$.
\end{proof}

\begin{computation}[a 4-number near-miss with only three common differences]
There is a 4-tuple
\[
(1040,2660,5520,8580)
\]
with
\[
\Dset(1040)\cap \Dset(2660)\cap \Dset(5520)\cap \Dset(8580)=\{32,67,256\}.
\]
In normalized square-translate form this becomes a $3$-shift/$5$-square configuration
\[
e=\{0,3465,64512\},\qquad u=\{32,72,108,152,188\}.
\]
\end{computation}

\begin{computation}[a $3\times 4$ configuration]
There is a triple
\[
(13608,29848,65968)
\]
with four common differences
\[
\Dset(13608)\cap \Dset(29848)\cap \Dset(65968)=\{18,282,477,1122\}.
\]
In normalized square-translate form this becomes a $4$-shift/$4$-square configuration
\[
e=\{0,79200,227205,1258560\},\qquad u=\{18,234,346,514\}.
\]
\end{computation}

\begin{computation}[search records]
The following exact search results were recorded in the notes.
\begin{enumerate}[label=\textup{(\roman*)}]
    \item No 4-tuple with four common differences was found in an exact streaming search up to $35{,}000$.
    \item No 4-tuple with four common differences was found in a direct search up to $50{,}000$.
    \item In a restricted search up to $10^6$, looking only for quartets with common differences at most $1000$, again no 4-tuple with four common differences was found.
\end{enumerate}
These are certified computations, not proofs of global nonexistence.
\end{computation}

\section{Anchored packets and divisor grids}

\subsection{The anchored formalism}

Fix an anchored common-difference set
\[
d_1<d_2<\cdots<d_r
\]
and write
\[
c_i:=d_i^2-d_1^2\qquad (2\le i\le r).
\]
For a supporting integer $N$, set
\[
x:=\sqrt{d_1^2+4N}.
\]
Then
\[
x^2+c_i=d_i^2+4N
\]
is a square for every $i$.

\begin{proposition}[anchored packet formulation]
A finite set $T=\{d_1,\dots,d_r\}$ has $m$ positive common supports
\[
N_1,\dots,N_m
\]
if and only if there exist positive integers
\[
x_1,\dots,x_m
\]
such that for every $i=2,\dots,r$ and every $t=1,\dots,m$,
\[
x_t^2+c_i \text{ is a square}.
\]
Equivalently, for each pair $(i,t)$ there exist positive integers $r_{i,t}$ such that
\[
c_i=r_{i,t}(2x_t+r_{i,t}).
\]
\end{proposition}

\begin{proof}
This is simply the identity
\[
x_t^2+c_i=(x_t+r_{i,t})^2
\iff c_i=r_{i,t}(2x_t+r_{i,t}).
\]
Conversely, from such a packet one recovers
\[
d_i=\sqrt{d_1^2+c_i},\qquad N_t=\frac{x_t^2-d_1^2}{4}.
\]
\end{proof}

\subsection{The exact cubic column law}

\begin{theorem}[column law]
For a fixed column $t$, the anchored divisor-grid variables satisfy
\[
c_j r_{i,t}-c_i r_{j,t}=r_{i,t}r_{j,t}(r_{j,t}-r_{i,t})
\qquad (i<j).
\]
Equivalently, if one sets
\[
y_{i,t}:=x_t+r_{i,t}=\sqrt{x_t^2+c_i},
\qquad
\delta_{ij,t}:=y_{j,t}-y_{i,t}=r_{j,t}-r_{i,t}>0,
\]
and
\[
\Delta_{ij}:=c_j-c_i,
\]
then
\[
\delta_{ik,t}=\delta_{ij,t}+\delta_{jk,t},
\]
and
\[
\Delta_{jk}\,\delta_{ij,t}-\Delta_{ij}\,\delta_{jk,t}
=
\delta_{ij,t}\delta_{jk,t}\delta_{ik,t}
\qquad (i<j<k).
\]
\end{theorem}

\begin{proof}
From
\[
c_i=r_{i,t}(2x_t+r_{i,t}),\qquad c_j=r_{j,t}(2x_t+r_{j,t}),
\]
subtracting yields
\[
c_j-c_i=2x_t(r_{j,t}-r_{i,t})+(r_{j,t}^2-r_{i,t}^2)
=(r_{j,t}-r_{i,t})(2x_t+r_{i,t}+r_{j,t}),
\]
from which
\[
2x_t=\frac{c_j-c_i}{r_{j,t}-r_{i,t}}-(r_{i,t}+r_{j,t}).
\]
Eliminating $x_t$ between the equations for $c_i$ and $c_j$ gives
\[
c_j r_{i,t}-c_i r_{j,t}=r_{i,t}r_{j,t}(r_{j,t}-r_{i,t}).
\]
Now write
\[
\Delta_{ij}=(y_{j,t}-y_{i,t})(y_{j,t}+y_{i,t})=\delta_{ij,t}(2x_t+r_{i,t}+r_{j,t}),
\]
which yields the displayed cocycle law after eliminating $x_t$.
\end{proof}

\begin{remark}
The important structural point is that the exact packet law is \emph{cubic}.  The correct object is a divisor-valued additive cocycle on row pairs, not a separable, bilinear, or Möbius law.
\end{remark}

\subsection{Completed packets of width $4$}

\begin{computation}[completed width-$4$ packets]
The following exact packets occur.
\begin{enumerate}[label=\textup{(\roman*)}]
    \item
    \[
    (c_2,c_3)=(2880,12285),\qquad
    x=(6,22,62,114),
    \]
    with
    \[
    r=
    \begin{pmatrix}
    48&36&20&12\\
    105&91&65&45
    \end{pmatrix}.
    \]
    \item
    \[
    (c_2,c_3)=(1365,3360),\qquad
    x=(2,22,46,134),
    \]
    with
    \[
    r=
    \begin{pmatrix}
    35&21&13&5\\
    56&40&28&12
    \end{pmatrix}.
    \]
    \item
    \[
    (c_2,c_3,c_4)=(79200,227205,1258560),\qquad
    x=(18,234,346,514),
    \]
    with
    \[
    r=
    \begin{pmatrix}
    264&132&100&72\\
    459&297&243&187\\
    1104&912&828&720
    \end{pmatrix}.
    \]
\end{enumerate}
In particular, the originally displayed $3$-column fragments extend to honest width-$4$ packets.
\end{computation}

\begin{proof}[Verification]
Each entry is checked directly from the identity
\[
c_i=r_{i,t}(2x_t+r_{i,t}).
\]
For example, in the $3\times 4$ packet,
\[
79200=264\cdot 300=132\cdot 600=100\cdot 792=72\cdot 1100,
\]
\[
227205=459\cdot 495=297\cdot 765=243\cdot 935=187\cdot 1215,
\]
\[
1258560=1104\cdot 1140=912\cdot 1380=828\cdot 1520=720\cdot 1748.
\]
\end{proof}

\subsection{The Möbius / projective-rank-$2$ obstruction}

\begin{theorem}[cross-ratio test on completed packets]
If a completed packet comes from a common Möbius parameter
\[
r_{i,t}=M_i(z_t)
\qquad
\left(
M_i(z)=\frac{\alpha_i z+\gamma_i}{\varepsilon_i z+\phi_i}
\right),
\]
then for any four columns the row cross-ratios must agree:
\[
\CR(r_{i,1},r_{i,2},r_{i,3},r_{i,4})
=
\CR(z_1,z_2,z_3,z_4).
\]
For the completed $3\times 4$ packet one has
\[
\CR(264,132,100,72)=\frac{205}{128},
\]
\[
\CR(459,297,243,187)=\frac{55}{34},
\]
\[
\CR(1104,912,828,720)=\frac{23}{14},
\]
which are distinct. Hence the packet is not of Möbius type. The two $2\times 4$ packets fail the same test.
\end{theorem}

\begin{proof}
The cross-ratio of four values of a Möbius function depends only on the underlying four parameter values. The explicit computations above are exact.
\end{proof}

\begin{theorem}[projective-rank-$2$ obstruction for width $5$]
Let $\kappa_i$ be distinct positive row constants. Assume
\[
r_{i,t}=M_i(z_t)=\frac{a_i z_t+b_i}{c_i z_t+d_i}
\]
for five distinct column parameters $z_t$, and
\[
\kappa_i=r_{i,t}(u_t+r_{i,t})
\]
for all $i,t$. Then such a width-$5$ divisor grid is impossible.
\end{theorem}

\begin{proof}
Define
\[
h_\kappa(w):=\frac{\kappa}{w}-w.
\]
Then
\[
u_t=h_{\kappa_i}(M_i(z_t))
\]
for every $i,t$. Hence for each pair $i,j$,
\[
h_{\kappa_i}(M_i(z_t))=h_{\kappa_j}(M_j(z_t))
\]
at the five distinct values $z_t$.

Now
\[
h_{\kappa_i}(M_i(z))
=
\frac{\kappa_i(c_i z+d_i)^2-(a_i z+b_i)^2}
{(a_i z+b_i)(c_i z+d_i)}
\]
is a rational function whose numerator after clearing denominators has degree at most $4$. Therefore equality at five distinct points forces an identity
\[
h_{\kappa_i}\circ M_i \equiv h_{\kappa_j}\circ M_j.
\]
Writing
\[
\psi:=M_j\circ M_i^{-1},
\]
this becomes
\[
h_{\kappa_j}\circ \psi = h_{\kappa_i}.
\]
Write
\[
\psi(z)=\frac{az+b}{cz+d},\qquad ad-bc\ne 0.
\]
Clearing denominators in
\[
\frac{\kappa_j}{\psi(z)}-\psi(z)=\frac{\kappa_i}{z}-z
\]
gives the quartic identity
\[
ac z^4+(-a^2+ad+bc+c^2\kappa_j)z^3+(-2ab-ac\kappa_i+bd+2cd\kappa_j)z^2
\]
\[
+(-ad\kappa_i-b^2-bc\kappa_i+d^2\kappa_j)z-bd\kappa_i=0.
\]
Thus
\[
ac=0,\qquad bd=0.
\]
Since $ad-bc\ne 0$, there are only two cases.

If $c=0$, then $b=0$, so $\psi(z)=\lambda z$. Substituting back into the identity forces $\lambda=1$ and hence $\kappa_i=\kappa_j$.

If $a=0$, then $d=0$, so $\psi(z)=\lambda/z$. Substituting back forces $\lambda=-\kappa_i=-\kappa_j$, again implying $\kappa_i=\kappa_j$.

This contradicts the distinctness of the row constants.
\end{proof}

\begin{corollary}
No divisor grid of width $5$ can come from the full projective-rank-$2$ family
\[
r_{i,t}
=
\frac{\alpha_i\beta_t+\gamma_i\delta_t}
{\varepsilon_i\beta_t+\phi_i\delta_t}.
\]
In particular, every common-parameter rational one-variable law for width $5$ is ruled out.
\end{corollary}

\subsection{Elliptic transfer for $3$ rows and genus-$5$ transfer for $4$ rows}

\begin{theorem}[exact $3$-row transfer law]
Fix
\[
0<c_2<c_3<c_4,
\qquad
B:=c_3-c_2,\qquad A:=c_4-c_3.
\]
For a column define
\[
u_t:=r_{3,t}-r_{2,t},\qquad
v_t:=r_{4,t}-r_{3,t}.
\]
Then the exact adjacent law is
\[
A u_t-B v_t=u_t v_t(u_t+v_t).
\]
Under the change of variables
\[
X_t=\frac{A u_t}{v_t},\qquad
Y_t=\frac{A u_t(u_t+v_t)}{v_t},
\]
the column law becomes the elliptic curve
\[
E_{A,B}:\quad Y^2=X(X-B)(X+A).
\]
Conversely,
\[
u=\frac{Y}{X+A},\qquad
v=\frac{A\,Y}{X(X+A)},\qquad
u+v=\frac{Y}{X}
\]
recover the column differences from a point on $E_{A,B}$.
\end{theorem}

\begin{proof}
From
\[
A u-B v=u v(u+v)
\]
one gets
\[
u^2\left(\frac{A u}{v}+A\right)
=
\frac{A u}{v}\left(\frac{A u}{v}-B\right).
\]
Substituting the definitions of $X$ and $Y$ gives
\[
Y^2=X(X-B)(X+A).
\]
The inverse formulas follow by direct substitution.
\end{proof}

\begin{theorem}[exact $4$-row transfer law]
Fix
\[
0<c_2<c_3<c_4<c_5,
\qquad
B=c_3-c_2,\quad A=c_4-c_3,\quad C=c_5-c_4.
\]
For a column define
\[
u=r_3-r_2,\qquad v=r_4-r_3,\qquad w=r_5-r_4.
\]
Then
\[
A u-B v=u v(u+v),
\qquad
C v-A w=v w(v+w).
\]
Solving for $u$ and $w$ gives
\[
u=\frac{A-v^2\pm Y_-}{2v},
\qquad
Y_-^2=v^4-2(A+2B)v^2+A^2,
\]
\[
w=\frac{-A-v^2\pm Y_+}{2v},
\qquad
Y_+^2=v^4+2(A+2C)v^2+A^2.
\]
The normalization $X_{A,B,C}$ of the biquadratic cover
\[
Y_-^2=f_-(v),\qquad
Y_+^2=f_+(v)
\]
has genus $5$.
\end{theorem}

\begin{proof}
The quadratic formulas are obtained by solving the two adjacent laws as quadratics in $u$ and $w$.  Put
\[
f_-(v)=v^4-2(A+2B)v^2+A^2,\qquad
f_+(v)=v^4+2(A+2C)v^2+A^2.
\]
One has
\[
f_-(v)=\bigl(v^2-(\sqrt{A+B}+\sqrt B)^2\bigr)\bigl(v^2-(\sqrt{A+B}-\sqrt B)^2\bigr),
\]
\[
f_+(v)=\bigl(v^2+(\sqrt{A+C}+\sqrt C)^2\bigr)\bigl(v^2+(\sqrt{A+C}-\sqrt C)^2\bigr),
\]
so the branch sets are four distinct real points for $f_-$ and four distinct purely imaginary points for $f_+$, hence disjoint.  The cover $X_{A,B,C}\to \PP^1_v$ has degree $4$, ramified over exactly $8$ points, with total ramification $16$. By Riemann--Hurwitz,
\[
2g(X_{A,B,C})-2=4(-2)+16=8,
\]
so $g(X_{A,B,C})=5$.
\end{proof}

\begin{corollary}[genus lower bound for exact one-parameter $4$-row laws]
Let $\mathcal B$ be a smooth irreducible curve carrying meromorphic functions
\[
r_2,r_3,r_4,r_5,x\in \Q(\mathcal B)
\]
such that
\[
c_i=r_i(2x+r_i)\qquad (i=2,3,4,5)
\]
identically on $\mathcal B$. Then $\mathcal B$ admits a nonconstant morphism to a genus-$5$ curve. In particular,
\[
g(\mathcal B)\ge 5.
\]
\end{corollary}

\begin{proof}
The functions
\[
u=r_3-r_2,\qquad v=r_4-r_3,\qquad w=r_5-r_4
\]
satisfy the equations of Theorem 4.10 identically, hence define a nonconstant map $\mathcal B\to X_{A,B,C}$.  By Riemann--Hurwitz, $g(\mathcal B)\ge g(X_{A,B,C})=5$.
\end{proof}

\subsection{Two exact packet saturations}

\begin{computation}[the packet $(79200,227205,1258560)$ is column-saturated]
For the packet
\[
(c_2,c_3,c_4)=(79200,227205,1258560),
\qquad
(B,A)=(148005,1031355),
\]
the only admissible positive values of $u\mid B$ are
\[
u\in\{115,143,165,195\}.
\]
They produce exactly the four source columns
\[
(u,v)=
(115,533),\ (143,585),\ (165,615),\ (195,645),
\]
equivalently
\[
x\in\{18,234,346,514\}.
\]
Hence no positive integer $C$ can yield a $4\times 5$ positive divisor grid extending this packet.
\end{computation}

\begin{proof}[Verification]
For fixed $u\mid B$, the exact divisor sieve is
\[
D_-(u):=\left(\frac Bu-u\right)^2-4c_2,
\qquad
D_+(u):=\left(\frac Bu+u\right)^2+4A.
\]
Positivity requires both $D_-(u)$ and $D_+(u)$ to be squares, together with
\[
r_2=\frac{\frac Bu-u-\sqrt{D_-(u)}}{2}>0,
\qquad
v=\frac{-\left(\frac Bu+u\right)+\sqrt{D_+(u)}}{2}>0.
\]
Exact divisor enumeration over $u\mid 148005$ yields only the four admissible values above. Thus the positive integral source packet has exactly four columns.
\end{proof}

\begin{computation}[a new exact $3$-row/$5$-column packet]
The triple
\[
(c_2,c_3,c_4)=(756000,15971200,45130176)
\]
has exactly five common positive columns
\[
x=\{330,870,2445,4155,10482\},
\]
with
\[
\sqrt{x^2+c_2}\in\{930,1230,2595,4245,10518\},
\]
\[
\sqrt{x^2+c_3}\in\{4010,4090,4685,5765,11218\},
\]
\[
\sqrt{x^2+c_4}\in\{6726,6774,7149,7899,12450\}.
\]
It is governed by the elliptic packet law for
\[
(A,B)=(29158976,15215200),
\]
and the five associated integral points on
\[
E_{A,B}:\quad Y^2=X(X-B)(X+A)
\]
are
\[
(16567600,32008603200),\ 
(20769280,75890949120),\ 
(24733060,112634355240),
\]
\[
(31071040,172257845760),\ 
(33066880,191655636480).
\]
Over these five columns the only common positive rows are exactly
\[
c_2,\ c_3,\ c_4.
\]
Thus this packet is a genuine new width-$5$ source packet, but it does not upgrade by adding a fourth row.
\end{computation}

\begin{proof}[Verification]
Direct substitution confirms that all three expressions $x_t^2+c_i$ are squares. The divisor-sieve data are
\[
u\in\{700,1520,2090,2860,3080\},
\]
with corresponding positive data
\[
r_2\in\{36,90,150,360,600\},\qquad
v\in\{1232,2134,2464,2684,2716\}.
\]
Exact enumeration of common positive rows above these five columns yields
\[
\bigcap_{t=1}^5 \{y^2-x_t^2:y>x_t\}
=
\{756000,15971200,45130176\}.
\]
\end{proof}

\section{Anchored spectra, row complexes, and square translates}

\subsection{The exact $B$-set identity}

\begin{lemma}
For $\delta,z>0$,
\[
2z\in \Bset(\delta)
\iff
z^2+\delta \text{ is a square}
\iff
\delta=x(x+2z)\text{ for some }x>0.
\]
\end{lemma}

\begin{proof}
By definition, $2z\in \Bset(\delta)$ means
\[
\delta=gh,\qquad h-g=2z,\qquad g\equiv h\pmod 2.
\]
Writing
\[
g=r-z,\qquad h=r+z
\]
gives
\[
\delta=r^2-z^2,
\]
which is equivalent to $z^2+\delta=r^2$.
\end{proof}

\subsection{The exact pairwise theorem for $B(a)\cap B(b)$}

\begin{theorem}[pairwise theorem]
Let $0<a<b$. Then
\[
2z\in \Bset(a)\cap \Bset(b)
\]
if and only if there exist positive integers $p,q$ such that
\[
pq=b-a
\]
and
\[
z^2=\left(\frac{q-p}{2}\right)^2-a.
\]
Equivalently,
\[
\Bset(a)\cap \Bset(b)=
\left\{\,2z>0:\exists\,pq=b-a,\ 4z^2=p^2+q^2-2(a+b)\,\right\}.
\]
\end{theorem}

\begin{proof}
By Lemma 5.1, $2z\in \Bset(a)\cap \Bset(b)$ iff there exist $m<n$ with
\[
m^2=z^2+a,\qquad n^2=z^2+b.
\]
Subtracting gives
\[
n^2-m^2=b-a.
\]
Set
\[
p:=n-m,\qquad q:=n+m.
\]
Then $pq=b-a$ and
\[
m=\frac{q-p}{2},
\]
hence
\[
z^2=m^2-a=\left(\frac{q-p}{2}\right)^2-a.
\]
The converse is immediate.
\end{proof}

\begin{corollary}
For fixed $a<b$, the set $\Bset(a)\cap \Bset(b)$ is finite, with a divisor-type bound controlled by the factorizations of $b-a$.
\end{corollary}

\subsection{Pairwise cap $4$ is false}

\begin{computation}[a pairwise spectrum of size $5$ inside a genuine same-parity sextuple]
Take
\[
(\Delta_2,\Delta_3)=(138240,760320),
\]
and
\[
s=(744,912,1104,1808,2928,6932).
\]
All six values lie in the same middle-interface class:
\[
\sqrt{s_t^2-4\Delta_2}=(24,528,816,1648,2832,6892),
\]
\[
\sqrt{s_t^2+4\Delta_3}=(1896,1968,2064,2512,3408,7148).
\]
Anchoring at $s_0=744$, the first two positive rows are
\[
\delta_1=\frac{912^2-744^2}{4}=69552,\qquad
\delta_2=\frac{1104^2-744^2}{4}=166320,
\]
and
\[
\Bset(69552)\cap \Bset(166320)=\{24,366,536,744,1896\}.
\]
Thus pairwise common spectrum can already have size $5$.
\end{computation}

\subsection{The exact anchored spectrum theorem}

\begin{theorem}[column-spectrum theorem]
Fix a $4$-column row complex with seed row
\[
R_0=(d_1,d_2,d_3,d_4),
\]
and positive rows
\[
R_t=(r_{t,1},r_{t,2},r_{t,3},r_{t,4})
\qquad (t=1,\dots,m),
\]
with
\[
r_{t,j}^2-d_j^2=4N_t.
\]
Write
\[
\delta_t:=4N_t,
\qquad
\Sigma:=\bigcap_{t=1}^m \Bset(\delta_t).
\]
Then $\Sigma$ is exactly the set of all doubled seed entries $2d$ compatible with all rows.  In particular,
\[
2d_j\in \Sigma\qquad (j=1,2,3,4),
\]
and a right extension exists if and only if $\Sigma$ contains an element strictly larger than $2d_4$.
\end{theorem}

\begin{proof}
For each seed entry $d_j$ one has
\[
\delta_t=(r_{t,j}-d_j)(r_{t,j}+d_j),
\]
and both factors have the same parity. Their difference is $2d_j$, so $2d_j\in \Bset(\delta_t)$ for every $t$, hence $2d_j\in \Sigma$.

Conversely, if $c\in \Sigma$, write $c=2d_*$. For each $t$ there exist $g_t,h_t>0$ with
\[
g_t h_t=\delta_t,\qquad h_t-g_t=c,\qquad g_t\equiv h_t\pmod 2.
\]
Set
\[
r_{t,*}:=\frac{g_t+h_t}{2}.
\]
Then
\[
r_{t,*}^2-d_*^2
=
\frac{(g_t+h_t)^2-(h_t-g_t)^2}{4}
=
g_t h_t
=
\delta_t.
\]
Thus $d_*$ is another seed entry compatible with all rows.
\end{proof}

\begin{theorem}[boundary embeddings in stitched form]
Fix a same-$\bmod 4$ set
\[
s_0<s_1<\cdots<s_m\subset \Mset(\Delta_2,\Delta_3),
\qquad s_t\equiv s_0\pmod 4.
\]
Set
\[
u_0:=\sqrt{s_0^2-4\Delta_2},\qquad
v_0:=\sqrt{s_0^2+4\Delta_3},
\]
and
\[
\delta_t:=\frac{s_t^2-s_0^2}{4},
\qquad
\Sigma(s_0;s_1,\dots,s_m):=\bigcap_{t=1}^m \Bset(\delta_t).
\]
Then
\[
u_0,\ s_0,\ v_0\in \Sigma(s_0;s_1,\dots,s_m).
\]
Moreover:
\begin{enumerate}[label=\textup{(\roman*)}]
    \item the left boundary survives if and only if there exists
    \[
    c_1\in \Sigma(s_0;s_1,\dots,s_m)\quad\text{with}\quad c_1<u_0,
    \]
    in which case
    \[
    \Delta_1=\frac{u_0^2-c_1^2}{4};
    \]
    \item the right boundary survives if and only if there exists
    \[
    c_5\in \Sigma(s_0;s_1,\dots,s_m)\quad\text{with}\quad c_5>v_0,
    \]
    in which case
    \[
    \Delta_4=\frac{c_5^2-v_0^2}{4}.
    \]
\end{enumerate}
Hence a $5$-column witness exists exactly when
\[
\Sigma(s_0;s_1,\dots,s_m)
\]
contains five ordered values
\[
c_1<u_0<s_0<v_0<c_5.
\]
\end{theorem}

\begin{proof}
The inclusions $u_0,s_0,v_0\in \Sigma$ are obtained by taking the factorizations coming from the left edge, the middle interface, and the right edge of the anchored row.  The boundary conditions are then immediate from the column-spectrum theorem.
\end{proof}

\begin{theorem}[anchored spectral-cap obstruction]
Fix $(\Delta_2,\Delta_3)$. Assume that for every candidate seed middle interface
\[
s_0\in \Mset(\Delta_2,\Delta_3),
\]
and every two larger interfaces
\[
s_i,s_j\in \Mset(\Delta_2,\Delta_3),\qquad
s_i\equiv s_j\equiv s_0\pmod 4,
\]
one has
\[
\left|
\Bset\!\left(\frac{s_i^2-s_0^2}{4}\right)
\cap
\Bset\!\left(\frac{s_j^2-s_0^2}{4}\right)
\right|
\le 4.
\]
Then no $4$-column seed based on $(\Delta_2,\Delta_3)$ with five positive rows can extend to a fifth column.
\end{theorem}

\begin{proof}
If such an extension existed, then for the chosen seed interface $s_0$ and the five positive interfaces $s_1,\dots,s_5$ the anchored spectrum
\[
\Sigma(s_0;s_1,\dots,s_5)
\]
would contain
\[
2d_1,\ 2d_2,\ 2d_3,\ 2d_4,\ 2d_5,
\]
hence at least five elements. But $\Sigma$ is contained in every pairwise intersection on the right-hand side, contradiction.
\end{proof}

\begin{corollary}[mod $4$ class obstruction]
A $k=5$ witness requires some $\Mset(\Delta_2,\Delta_3)$ to contain at least six elements in a single residue class modulo $4$.
\end{corollary}

\subsection{Applications}

\begin{computation}[the verified $4$-column seed is spectrally capped at $4$]
For the seed
\[
(18,282,477,1122)
\]
with positive rows
\[
(234,366,531,1146),\qquad
(346,446,589,1174),\qquad
(514,586,701,1234),
\]
the positive-row parameters are
\[
\delta_1=54432,\qquad
\delta_2=119392,\qquad
\delta_3=263872.
\]
One has
\[
\Bset(54432)\cap \Bset(119392)=\{36,564,954,2244\},
\]
and
\[
\Bset(54432)\cap \Bset(119392)\cap \Bset(263872)=\{36,564,954,2244\}.
\]
So the anchored spectrum is already capped at $4$.
\end{computation}

\begin{computation}[the $7$-interface near-miss is structurally dead]
For
\[
(\Delta_2,\Delta_3)=(48960,196560),
\]
one has
\[
\Mset(\Delta_2,\Delta_3)=\{458,496,528,606,1264,2738,8166\}.
\]
It splits mod $4$ as
\[
\{458,606,2738,8166\}\qquad\text{and}\qquad \{496,528,1264\}.
\]
In the first class, with $s_0=458$,
\[
\Bset\!\left(\frac{606^2-458^2}{4}\right)\cap
\Bset\!\left(\frac{2738^2-458^2}{4}\right)
=
\{118,458,998\}.
\]
In the second class, with $s_0=496$,
\[
\Bset\!\left(\frac{528^2-496^2}{4}\right)\cap
\Bset\!\left(\frac{1264^2-496^2}{4}\right)
=
\{224,496,1016\}.
\]
Thus neither residue class even supports a $4$-column seed, let alone a $5$-column extension.
\end{computation}

\subsection{Quartic and hyperelliptic reformulations}

\begin{proposition}[anchored quartic/hyperelliptic reformulation]
Let
\[
y_0<y_1<\cdots<y_5
\]
be same-parity integers. A $k=5$ anchored witness is equivalent to the existence of positive integers
\[
\Delta_1>\Delta_2>0,\qquad \Delta_4>\Delta_3>0
\]
such that
\[
y_t^2-\Delta_1,\quad
y_t^2-\Delta_2,\quad
y_t^2+\Delta_3,\quad
y_t^2+\Delta_4
\]
are all squares for every $t=0,\dots,5$.
Equivalently, the six numbers
\[
T_t:=y_t^2
\]
are six square integral abscissae on the quartic
\[
E:\quad W^2=(T-\Delta_1)(T-\Delta_2)(T+\Delta_3)(T+\Delta_4),
\]
a genus-$1$ curve in $T$.  Equivalently again, the twelve points $\pm y_t$ lie on the genus-$3$ hyperelliptic curve
\[
H:\quad W^2=(Y^2-\Delta_1)(Y^2-\Delta_2)(Y^2+\Delta_3)(Y^2+\Delta_4).
\]
\end{proposition}

\begin{proof}
This is the exact fixed-shift reformulation
\[
X+y_t^2\in \square \quad (t=0,\dots,5)
\]
for two extra left levels and two extra right levels.  Writing $T=z^2$ for a spectrum value yields the quartic equation.
\end{proof}

\section{Canonical divisor-seed dynamics and row complexes}

\subsection{The exact recursion}

\begin{definition}
For a finite set $T\subset \Z_{>0}$ with $\Sigma(T)\neq\emptyset$, define
\[
H(T):=\min \Sigma(T).
\]
For each $d\in T$, write
\[
H(T)=u_d(u_d+d),
\]
and define
\[
S(T):=\{\,2u_d+d : d\in T\,\}
=
\{\sqrt{d^2+4H(T)}:d\in T\}.
\]
\end{definition}

\begin{theorem}[exact sum-difference duality]
For every finite set $T$ with $\Sigma(T)\neq\emptyset$,
\[
\Sigma(T)=\{H(T)\}\ \sqcup\ \bigl(H(T)+\Sigma(S(T))\bigr).
\]
\end{theorem}

\begin{proof}
Let $N>H(T)$ and $d\in T$. Writing
\[
H(T)=u_d(u_d+d),
\]
and
\[
4N+d^2=x_d^2,
\]
one gets
\[
x_d=2t_d+2u_d+d
\]
for some $t_d\ge 0$, hence
\[
N-H(T)=t_d(t_d+2u_d+d).
\]
Thus $2u_d+d\in \Dset(N-H(T))$, so $N-H(T)\in \Sigma(S(T))$.

Conversely, if $M\in \Sigma(S(T))$, then
\[
M=t_d(t_d+2u_d+d)
\]
for every $d\in T$. Hence
\[
H(T)+M=(u_d+t_d)(u_d+d+t_d),
\]
so $d\in \Dset(H(T)+M)$ for every $d\in T$.
\end{proof}

\subsection{Exact decoding of the main examples}

\begin{example}[the pair $\{1,11\}$]
For
\[
T_0=\{1,11\},
\qquad
H(T_0)=12,
\]
one finds
\[
S(T_0)=\{7,13\},\qquad H(\{7,13\})=30,\qquad S(\{7,13\})=\{13,17\},
\]
\[
H(\{13,17\})=168,\qquad S(\{13,17\})=\{29,31\},
\]
and $\Sigma(\{29,31\})=\varnothing$. Therefore
\[
\Sigma(\{1,11\})=\{12,\ 12+30,\ 12+30+168\}=\{12,42,210\}.
\]
\end{example}

\begin{example}[the triple $\{6,54,111\}$]
For
\[
T_0=\{6,54,111\},
\qquad H(T_0)=112,
\]
one gets
\[
S(T_0)=\{22,58,113\},
\qquad H=840,
\qquad S=\{62,82,127\},
\]
\[
H=2288,
\qquad S=\{114,126,159\},
\]
and the chain stops. Hence
\[
\Sigma(\{6,54,111\})=\{112,952,3240\}.
\]
\end{example}

\begin{example}[the triple $\{2,37,58\}$]
For
\[
T_0=\{2,37,58\},
\qquad H(T_0)=120,
\]
one gets the chain
\[
\{2,37,58\}\to \{22,43,62\}\to \{46,59,74\}\to \{134,139,146\},
\]
with
\[
H=120,\ 408,\ 3960.
\]
Thus
\[
\Sigma(\{2,37,58\})=\{120,528,4488\}.
\]
\end{example}

\begin{example}[the four-difference packet]
For
\[
T_0=\{18,282,477,1122\},
\qquad H(T_0)=13608,
\]
one gets
\[
S(T_0)=\{234,366,531,1146\},\qquad H=16240,
\]
\[
S=\{346,446,589,1174\},\qquad H=36120,
\]
\[
S=\{514,586,701,1234\},
\]
and the chain stops. Hence
\[
\Sigma(\{18,282,477,1122\})=\{13608,29848,65968\}.
\]
\end{example}

\subsection{Local dynamics and the row-complex theorem}

\begin{proposition}[pair-gap contraction]
At stage $r$, write
\[
H_r=H(T_r),
\qquad
H_r=u_{r,j}(u_{r,j}+d_{r,j}),
\qquad
m_{r,j}:=2u_{r,j}+d_{r,j}=d_{r+1,j}.
\]
Then for every pair $i<j$,
\[
d_{r+1,j}-d_{r+1,i}
=
(d_{r,j}-d_{r,i})
\frac{H_r-u_{r,i}u_{r,j}}{H_r+u_{r,i}u_{r,j}}.
\]
\end{proposition}

\begin{proof}
From
\[
d_{r,j}=\frac{H_r}{u_{r,j}}-u_{r,j},
\qquad
d_{r+1,j}=\frac{H_r}{u_{r,j}}+u_{r,j},
\]
subtract the formulas for $i$ and $j$ and simplify.
\end{proof}

\begin{theorem}[row-complex theorem]
Fix an initial ordered seed
\[
T_0=(d_1<\cdots<d_k),
\]
and set
\[
\Delta_{ij}:=d_j^2-d_i^2.
\]
Let $\mathcal R(T_0)$ be the set of all ordered $k$-tuples
\[
R=(r_1<\cdots<r_k)
\]
such that
\begin{enumerate}[label=\textup{(\roman*)}]
    \item $r_j\equiv d_j\pmod 2$ for all $j$;
    \item $r_j^2-r_i^2=\Delta_{ij}$ for all $i<j$.
\end{enumerate}
Then the map
\[
R\longmapsto N(R):=\frac{r_j^2-d_j^2}{4}
\]
is a bijection
\[
\mathcal R(T_0)\longleftrightarrow \Sigma(T_0)\sqcup\{0\}.
\]
In particular, the canonical chain is exactly the ordered list of rows in $\mathcal R(T_0)$.
\end{theorem}

\begin{proof}
If $R$ is a row in the canonical chain, then
\[
r_j^2-d_j^2=4N
\]
is independent of $j$, so $R\in \mathcal R(T_0)$ and $N(R)=N$.
Conversely, if $R\in \mathcal R(T_0)$, then
\[
N:=\frac{r_j^2-d_j^2}{4}
\]
is a positive integer independent of $j$, so $d_j\in \Dset(N)$ for every $j$. Thus $N\in \Sigma(T_0)$. The case $N=0$ corresponds to the seed row itself.
\end{proof}

\begin{proposition}[stitched factor-pair description of a row]
For adjacent columns define
\[
\Delta_i:=d_{i+1}^2-d_i^2,\qquad
g_i:=r_{i+1}-r_i,\qquad
h_i:=r_{i+1}+r_i.
\]
Then every row satisfies
\[
g_i h_i=\Delta_i,
\qquad
h_i+g_i=h_{i+1}-g_{i+1}.
\]
Conversely, any stitched chain of factor pairs satisfying these relations reconstructs a row.
\end{proposition}

\begin{proof}
The identities are immediate from
\[
r_{i+1}^2-r_i^2=\Delta_i,
\qquad
2r_{i+1}=h_i+g_i.
\]
Conversely,
\[
r_1=\frac{h_1-g_1}{2},
\qquad
r_{i+1}=\frac{h_i+g_i}{2}
\]
recovers the row.
\end{proof}

\subsection{The terminal stitch criterion}

\begin{theorem}[terminal stitch criterion]
Let
\[
R_t=(r_{t,1},\dots,r_{t,k})
\qquad (t=0,\dots,s)
\]
be a full $k$-column row complex. Set
\[
s_t:=2r_{t,k}.
\]
Then the complex extends to a $(k+1)$-st column if and only if there exists $\Lambda>0$ such that
\[
s_t^2+4\Lambda
\]
is a square for every $t$.  Equivalently,
\[
\Lambda\in \bigcap_t \{x(x+s_t):x\in \Z_{>0}\}.
\]
\end{theorem}

\begin{proof}
To extend by a final gap $\Lambda$, one needs numbers $e_t$ with
\[
e_t^2-r_{t,k}^2=\Lambda.
\]
Set
\[
u_t=e_t-r_{t,k},\qquad v_t=e_t+r_{t,k}.
\]
Then
\[
u_t v_t=\Lambda,\qquad v_t-u_t=2r_{t,k}=s_t.
\]
Thus
\[
s_t^2+4\Lambda=(v_t-u_t)^2+4u_tv_t=(u_t+v_t)^2
\]
is a square. The converse is the same argument in reverse.
\end{proof}

\begin{example}[exact reproduction of the main chains]
For the $3$-column seed $(6,54,111)$, the last-column sequence
\[
111,\ 113,\ 127,\ 159
\]
has no common square translate, so the row complex does not extend to $4$ columns.

For the first three columns of the seed $(18,282,477,1122)$, the last-column sequence
\[
477,\ 531,\ 589,\ 701
\]
has the unique common square translate
\[
\Lambda=1031355,
\]
producing
\[
1122,\ 1146,\ 1174,\ 1234.
\]
The full $4$-column row complex then fails to extend because
\[
1122,\ 1146,\ 1174,\ 1234
\]
have no common square translate.
\end{example}

\begin{computation}[long chains exist for $k=3$ but die at the next column]
The $3$-column seed
\[
(126,894,1986)
\]
has row complex of length $8$:
\[
(126,894,1986),
\ (236,916,1996),
\ (448,992,2032),
\ (576,1056,2064),
\]
\[
(828,1212,2148),
\ (2368,2528,3088),
\ (5404,5476,5756),
\ (16308,16332,16428).
\]
Thus
\[
|\Sigma(T_0)|=7.
\]
However, the last-column sequence
\[
1986,\ 1996,\ 2032,\ 2064,\ 2148,\ 3088,\ 5756,\ 16428
\]
has no common square translate, so the complex does not extend to four columns.
\end{computation}

\begin{computation}[exact searches in the row-complex model]
The notes recorded the following exact search outcomes.
\begin{enumerate}[label=\textup{(\roman*)}]
    \item Among canonical $4$-element seeds with $H(T_0)\le 20000$, one has
    \[
    \max |\Sigma(T_0)|=3.
    \]
    None of the resulting $4$-row complexes extends to a $5$th column.
    \item Among canonical $5$-element seeds with $H(T_0)\le 16000$, one has
    \[
    \max |\Sigma(T_0)|=2.
    \]
    A sample maximizer is
    \[
    T_0=(2,43,97,178,362),
    \]
    with row complex
    \[
    (2,43,97,178,362)\to (54,69,111,186,366)\to (306,309,321,354,474).
    \]
\end{enumerate}
These are exact search records, not general theorems.
\end{computation}

\section{Static divisor-seed geometry}

\begin{proposition}[pairwise factorization formula]
Fix a seed $H$ and divisors $u_j\mid H$, and define
\[
d_j=\frac{H}{u_j}-u_j.
\]
If
\[
c_{ij}:=\frac{H}{u_i u_j}\in \Z,
\]
then
\[
d_i^2-d_j^2
=
(u_j^2-u_i^2)(c_{ij}^2-1)
=
(u_j-u_i)(u_j+u_i)(c_{ij}-1)(c_{ij}+1).
\]
\end{proposition}

\begin{proof}
Expand
\[
d_j^2=\left(\frac{H}{u_j}-u_j\right)^2
\]
and simplify.
\end{proof}

\begin{proposition}[the static support curve]
Fix a seed
\[
T=\{d_1,\dots,d_r\}\subseteq \Dset(H),
\]
and set
\[
m_j:=\sqrt{d_j^2+4H}.
\]
The extra supports $N>H$ correspond exactly to integral points on the projective curve
\[
C_T:\quad x_j^2-x_1^2=(m_j^2-m_1^2)z^2\qquad (2\le j\le r).
\]
If $C_T$ is smooth, then $C_T$ is a complete intersection of $r-1$ quadrics in $\PP^r$ and has genus
\[
g(C_T)=1+2^{r-2}(r-3).
\]
\end{proposition}

\begin{proof}
The equations
\[
x_j^2=m_j^2+4M
\]
for a new support $H+M$ eliminate $M$ and give the displayed complete intersection.  Its degree is $2^{r-1}$.  By adjunction, for a smooth complete intersection of $r-1$ quadrics in $\PP^r$ one has
\[
K_{C_T}=(r-3)H|_{C_T},
\]
hence
\[
2g(C_T)-2=(r-3)2^{r-1}.
\]
\end{proof}

\begin{corollary}[high-genus obstruction for the static seed approach]
For $r=2$ the static support curve is genus $0$; for $r=3$ it is genus $1$; for $r=4$ it is genus $5$; and for all $r\ge 4$ it is high genus. In particular, pairwise factorability of the numbers $d_i^2-d_j^2$ cannot, by itself, force a uniform static construction for arbitrary $r$.
\end{corollary}

\section{Square-translate geometry and the genus-$1$/genus-$5$ jump}

\subsection{Normalized square-translate form}

\begin{proposition}[normalized square-translate equivalence]
The original $k\times k$ problem is equivalent to the following:
find shifts
\[
e_1=0,e_2,\dots,e_k
\]
and squares
\[
u_0^2,u_1^2,\dots,u_k^2
\]
such that every
\[
u_j^2+e_i
\]
is a square.  Equivalently, if one puts
\[
x_j:=u_j^2-u_0^2,\qquad d_i^2:=u_0^2+e_i,
\]
then
\[
x_j+d_i^2
\]
is a square for all $i,j$.
\end{proposition}

\subsection{The exact genus-$1$ model for three shifts}

\begin{theorem}[the $3$-shift curve is elliptic]
Fix shifts $(0,p,q)$. The common-square curve
\[
u^2+p=v^2,\qquad u^2+q=w^2
\]
is genus $1$.  A Weierstrass model is
\[
E_{p,q}:\quad Y^2=X(X+4p)(X+4q).
\]
The birational maps are
\[
X=4(u+v)(u+w),\qquad
Y=8(u+v)(u+w)(v+w),
\]
and
\[
u=\frac{16pq-X^2}{4Y},\qquad
v=\frac{X^2+8pX+16pq}{4Y},\qquad
w=\frac{X^2+8qX+16pq}{4Y}.
\]
\end{theorem}

\begin{proof}
Set
\[
w-v=\frac{q-p}{t},\qquad w+v=t.
\]
Then
\[
(2tu)^2=t^4-2(p+q)t^2+(q-p)^2,
\]
an even quartic, hence a genus-$1$ curve.  The displayed birational maps are checked directly by substitution.
\end{proof}

\begin{proposition}[a moving elliptic family]
Take
\[
p=(4t^2-1)^2,\qquad q=(t^2-4)^2.
\]
Then
\[
E_t:\quad Y^2=X\bigl(X+4(4t^2-1)^2\bigr)\bigl(X+4(t^2-4)^2\bigr)
\]
has the rational point
\[
P_t=\Bigl(4(2t+1)^2(t+2)^2,\ 40(2t+1)^2(t+2)^2(t^2+1)\Bigr).
\]
At $t=3$, the specialization
\[
E_3:\quad Y^2=X(X+4900)(X+100)
\]
has
\[
P_3=(4900,490000),
\]
and the notes record that $12P_3\ne O$, so $P_t$ is non-torsion in the family.  Consequently, for each $n\ge 1$ the point $[n]P_t$ yields a rational common-square solution for the three shifts $(0,p,q)$.
\end{proposition}

\subsection{The genus-$5$ jump for four shifts}

\begin{theorem}[the $4$-shift curve is genus $5$]
Fix shifts $(0,p,q,r)$. The common-square curve
\[
u^2+p=v^2,\qquad u^2+q=w^2,\qquad u^2+r=z^2
\]
is a smooth complete intersection of three quadrics in $\PP^4$, hence has genus
\[
g=5.
\]
Equivalently, after projecting to the elliptic quotient $E_{p,q}$, the curve is the fiber product
\[
C_{p,q,r}=E_{p,q}\times_{\PP^1_X} Q_{p,q,r},
\]
where
\[
Q_{p,q,r}:\quad Z^2=(16pq-X^2)^2+16r\,X(X+4p)(X+4q).
\]
\end{theorem}

\begin{proof}
The genus formula for a smooth complete intersection of three quadrics in $\PP^4$ is
\[
g=1+(3-2)2^2=5.
\]
To obtain the fiber-product description, use the inverse map from Theorem 7.2:
\[
u=\frac{16pq-X^2}{4Y}.
\]
Substituting into $u^2+r=z^2$ and clearing denominators yields the equation for $Q_{p,q,r}$.
\end{proof}

\begin{example}[normalized seeds]
The following examples from Section 3 become:
\[
(112,952,3240),\ \{6,54,111\}
\Longleftrightarrow
e=\{0,2880,12285\},\quad u=\{6,22,62,114\},
\]
\[
(13608,29848,65968),\ \{18,282,477,1122\}
\Longleftrightarrow
e=\{0,79200,227205,1258560\},\quad u=\{18,234,346,514\},
\]
\[
(1040,2660,5520,8580),\ \{32,67,256\}
\Longleftrightarrow
e=\{0,3465,64512\},\quad u=\{32,72,108,152,188\}.
\]
Thus the near-miss quartet is really a $3$-shift/$5$-square success, while the second seed is a $4$-shift/$4$-square success.
\end{example}

\section{The non-split seed-centered deformation}

\subsection{The basic seed geometry}

\begin{proposition}[the non-split seed]
Take
\[
(p,q,r)=(79200,227205,1258560).
\]
Then
\[
(r-p)(r-q)=1179360\cdot 1031355
=2^5\cdot 3^6\cdot 5^2\cdot 7\cdot 13^2\cdot 41\cdot 43,
\]
so the seed is genuinely non-split.

The associated genus-$2$ curve is
\[
C_{\mathrm{seed}}:\quad
y^2=(z^2-5034240)(z^2-4717440)(z^2-4125420).
\]
It carries the rational points
\[
(2244,19370016),\quad
(2292,363815712),\quad
(2348,727136992),\quad
(2468,1689152032),
\]
together with their $z\mapsto -z$ partners.
\end{proposition}

\begin{proof}
The non-split factorization is exact. The rational points come from the four rows of the seed packet:
\[
z=2\sqrt{u^2+r}\in\{2244,2292,2348,2468\}.
\]
Direct substitution verifies the coordinates.
\end{proof}

\begin{proposition}[the quotient elliptic curve and a divisor relation]
The quotient by the bielliptic involution $z\mapsto -z$ is
\[
E_{\mathrm{seed}}:\quad
Y^2=X(X+316800)(X+908820),
\qquad
X=z^2-4r=4u^2,\quad Y=y.
\]
The square-$X$ points are
\[
P_{18}=(1296,19370016),\quad
P_{234}=(219024,363815712),\quad
P_{346}=(478864,727136992),\quad
P_{514}=(1056784,1689152032).
\]
Moreover
\[
P_{18},\ P_{234},\ P_{514}
\]
are collinear on
\[
Y=1582X+17319744,
\]
so
\[
P_{18}+P_{234}+P_{514}=O
\]
on $E_{\mathrm{seed}}$. Pulling the line back to $C_{\mathrm{seed}}$ gives
\[
(1582z^2-7946847936)^2
-
(z^2-5034240)(z^2-4717440)(z^2-4125420)
\]
\[
=-(z^2-2244^2)(z^2-2292^2)(z^2-2468^2).
\]
\end{proposition}

\begin{proof}
The quotient equation follows from
\[
z^2-4717440=(z^2-4r)+316800,\qquad
z^2-4125420=(z^2-4r)+908820.
\]
Direct substitution shows that $P_{18},P_{234},P_{514}$ lie on the displayed line, and the factorization is an exact pullback of that line to the genus-$2$ curve.
\end{proof}

\begin{proposition}[cross-ratio obstruction to the old square-cross-ratio pencil]
For the seed core $\{0,79200,227205\}$, the normalized cross-ratio is
\[
\lambda_{\mathrm{seed}}=\frac{227205}{79200}=\frac{459}{160}.
\]
Its $S_3$-orbit is
\[
\frac{459}{160},\quad
\frac{160}{459},\quad
-\frac{299}{160},\quad
-\frac{160}{299},\quad
\frac{299}{459},\quad
\frac{459}{299},
\]
and no element of this orbit is a rational square.  The other three normalized $3$-shift cores from the same $4$-shift set give the cross-ratios
\[
\frac{874}{55},\qquad \frac{27968}{5049},\qquad \frac{2016}{253},
\]
and again none of their $S_3$-orbits contains a rational square. Therefore this seed does \emph{not} lie on the older square-cross-ratio pencil
\[
\lambda(t)=\left(\frac{t^2-4}{4t^2-1}\right)^2.
\]
\end{proposition}

\begin{proof}
The explicit orbit calculation is finite and exact.  A positive rational number is a rational square only if both numerator and denominator are squares in lowest terms, which fails for each value listed.
\end{proof}

\subsection{The seed-centered moving family}

\begin{proposition}[the $3$-point deformation base]
Keep
\[
r=1258560,\qquad a=4r=5034240,
\]
and keep the three square-$X$ values
\[
X_1=1296,\qquad X_2=219024,\qquad X_3=1056784
\]
fixed.  Let the quotient line move:
\[
Y=mX+n_0,\qquad n_0:=17319744.
\]
Then coefficient comparison in
\[
(mX+n_0)^2-X(X+4p)(X+4q)=-(X-1296)(X-219024)(X-1056784)
\]
gives
\[
4(p+q)=m^2-1277104,
\]
\[
16pq=233114505984+34639488\,m,
\]
and hence
\[
\Delta^2=(m-1532)(m-524)(m+596)(m+1460),
\]
where
\[
p=\frac{m^2-1277104-\Delta}{8},
\qquad
q=\frac{m^2-1277104+\Delta}{8}.
\]
Thus the seed-centered deformation base is the genus-$1$ curve
\[
\mathcal B_3:\quad
\Delta^2=(m-1532)(m-524)(m+596)(m+1460).
\]
\end{proposition}

\begin{proof}
The coefficient comparison is direct. The seed point $(m,\Delta)=(1582,592020)$ recovers
\[
(p,q)=(79200,227205).
\]
\end{proof}

\begin{proposition}[moving quotient curve and square-$X$ cover]
On the base $\mathcal B_3$, the moving quotient elliptic curve is
\[
E_{m,\Delta}:\quad Y^2=X(X+4p)(X+4q),
\]
and the moving genus-$2$ curve is
\[
C_{m,\Delta}:\quad
y^2=(z^2-4r)(z^2-4(r-p))(z^2-4(r-q)).
\]
The quadratic divisor section is
\[
y=mz^2+17319744-5034240m.
\]
On the square-$X$ cover
\[
\Sigma_{m,\Delta}:\quad W^2=(U^2+p)(U^2+q),
\]
the three built-in sections are
\[
R_{18}=(18,9m+120276),\qquad
R_{234}=(234,117m+9252),\qquad
R_{514}=(514,257m+4212).
\]
\end{proposition}

\begin{proof}
Substitute $X=4U^2$, $Y=8UW$ into the equation of $E_{m,\Delta}$.
\end{proof}

\subsection{The fixed $U_4=346$ slice}

\begin{proposition}[the genus-$3$ base with fixed fourth square-$X$ point]
Imposing the fourth square-$X$ point $U=346$ gives
\[
\mathcal B_4:\quad
\Delta^2=(m-1532)(m-524)(m+596)(m+1460),
\]
\[
v^2=29929m^2+2164968m-9320868336.
\]
The induced fifth square-$X$ section is
\[
U_5=\frac{1557m+9v+20807748}{59696}.
\]
At the seed point $(m,\Delta,v)=(1582,592020,262694)$, one gets
\[
U_5=\frac{3006}{7}.
\]
\end{proposition}

\begin{proof}
This is the quartic-interpolation computation on the square-$X$ cover through
\[
R_{234},\ R_{346},\ R_{514}.
\]
The specialization is exact.
\end{proof}

\begin{proposition}[simplified fixed-slice base and lift curve]
Eliminating $m,v$ in favor of
\[
x:=U_5
\]
gives
\[
m(x)=\frac{29848x^2-20807748x+3632717160}{1557x-541242},
\]
\[
v(x)=\frac{56(92209x^2-64107108x+11122790724)}{9(173x-60138)}.
\]
Substituting into the base equation yields the hyperelliptic genus-$3$ curve
\[
\mathcal B_4^{\mathrm{hyp}}:\quad Y^2=7P_8(x),
\]
where
\[
P_8(x)=
(13x-4698)(13x-4554)(13x-4338)
\]
\[
\times
(41x-17496)(41x-14886)(41x-14130)(41x-13626)
(91x-25146).
\]
The lift condition
\[
Z^2=4(r+x^2)=4(x^2+1258560)
\]
defines the curve
\[
\mathcal L:\quad
Y^2=7P_8(x),\qquad Z^2=4(x^2+1258560).
\]
The curve $\mathcal L$ has genus $7$.
\end{proposition}

\begin{proof}
The formulas for $m(x)$ and $v(x)$ follow by elimination.  The model for $\mathcal L$ is immediate.  Since $\mathcal B_4^{\mathrm{hyp}}$ is genus $3$ and the double cover to the $x$-line is branched at $16$ points in the fixed-slice model, the resulting lift curve has genus $7$.
\end{proof}

\begin{proposition}[quotients of the fixed lift curve]
The lift curve
\[
\mathcal L:\quad
Y^2=7P_8(x),\qquad Z^2=4(x^2+1258560)
\]
has the involutions
\[
\iota_Y:(x,Y,Z)\mapsto (x,-Y,Z),\qquad
\iota_Z:(x,Y,Z)\mapsto (x,Y,-Z).
\]
Their quotients are:
\[
\mathcal L/\langle \iota_Y\rangle:\ Z^2=4(x^2+1258560)
\quad\text{(genus $0$)},
\]
\[
\mathcal L/\langle \iota_Z\rangle:\ Y^2=7P_8(x)
\quad\text{(genus $3$)},
\]
\[
\mathcal L/\langle \iota_Y\iota_Z\rangle:\ W^2=7P_8(x)(x^2+1258560)
\quad\text{(genus $4$)}.
\]
Hence
\[
\Jac(\mathcal L)\sim \Jac(\mathcal B_4^{\mathrm{hyp}})\times \Jac(H_4),
\]
where
\[
H_4:\quad W^2=7P_8(x)(x^2+1258560).
\]
In the notes, the branch set of $H_4$ was checked to have trivial stabilizer, so no further low-genus quotient was found on this fixed slice.
\end{proposition}

\begin{remark}
The fixed slice $U_4=346$ is therefore very rigid: the obstruction is not a cheap local obstruction, but a genuinely genus-$7$ lift problem with no obvious hidden elliptic or genus-$2$ reduction.
\end{remark}

\subsection{The moving $(U_4,U_5)$ surface}

\begin{proposition}[general fifth square-$X$ section]
Let
\[
R_u=(u,w)\in \Sigma_{m,\Delta}
\]
be a moving fourth square-$X$ section. Then the induced fifth section is
\[
U_5(u,w)=\frac{u(9m+120276)+18w}{u^2-18^2}.
\]
Moreover, eliminating $\Delta$ from
\[
w^2=(u^2+p)(u^2+q)
\]
gives
\[
4w^2=u^2m^2+8659872m+4u^4-1277104u^2+58278626496.
\]
Eliminating $w$ between these equations yields a linear formula for $m$:
\[
m=
\frac{(u^2-324)U_5^2-240552\,uU_5-324u^2+14569656624}
{18(uU_5-120276)}.
\]
\end{proposition}

\begin{proof}
The formula for $U_5$ is the square-$X$ representative of the group-law difference $P_u-P_{18}$ on the quotient elliptic curve.  The elimination of $\Delta$ and then $w$ is direct symbolic algebra.
\end{proof}

\begin{proposition}[the two-parameter lift surface]
Impose the two lift conditions
\[
Z_4^2=4(r+u^2),\qquad Z_5^2=4(r+U_5^2).
\]
Parametrize the two conics by
\[
u=\frac{r-s^2}{2s},\qquad Z_4=\frac{r+s^2}{s},
\]
\[
U_5=\frac{r-t^2}{2t},\qquad Z_5=\frac{r+t^2}{t}.
\]
Then $m$ becomes an explicit rational function $m(s,t)$, and the only remaining condition is the base equation
\[
\Delta^2=(m(s,t)-1532)(m(s,t)-524)(m(s,t)+596)(m(s,t)+1460).
\]
Thus the natural seed-centered deformation is the two-parameter pullback surface in $(s,t)$, not the rigid fixed genus-$7$ fiber.
\end{proposition}

\begin{remark}
This moving surface is the main positive geometric lead remaining in the notes.  The fixed $U_4=346$ slice appears too rigid, but the moving two-parameter surface does not immediately collapse to a high-genus obstruction.
\end{remark}

\section{Literature map and recorded status}

\subsection{Direct status}

\begin{observation}[direct status recorded in the notes]
As of the status check recorded in the notes (April 2026), the full problem remained open.  The direct published chain is:
\begin{enumerate}[label=\textup{(\roman*)}]
    \item Erd\H{o}s--Rosenfeld (1997): the case $k=2$;
    \item Jim\'enez-Urroz (1999): the case $k=3$;
    \item Bremner (2019): the case $k=4$, including infinitely many $4$-tuples with four common factor differences.
\end{enumerate}
No later published proof of the full statement or of the case $k\ge 5$ was located in the notes.
\end{observation}

\subsection{Directly adjacent literature recorded in the notes}

The notes repeatedly identified the following neighboring literatures as most relevant.
\begin{enumerate}[label=\textup{(\alph*)}]
    \item \textbf{Lattice points on hyperbolas and close factorizations.}  This includes Cilleruelo--Jim\'enez-Urroz, Chan's papers on perfect squares, almost squares, and close factorizations, and later Pell-type refinements.
    \item \textbf{Square-translate and sumset-in-squares problems.}  The problem is a special case of asking for finite sets $X,Y$ with $X+Y$ contained in the squares.
    \item \textbf{Simultaneous Pell and elliptic-curve mechanisms.}  This is the line most closely aligned with the $k=4$ theorem and with the geometric reductions found in these notes.
    \item \textbf{Short-interval divisor phenomena.}  These provide structural obstructions and heuristics, especially around divisors near $\sqrt n$.
\end{enumerate}

\subsection{Recorded attack surfaces}

\begin{observation}[what the notes say to stop doing]
The following naive routes are decisively ruled out or strongly disfavored by the notes.
\begin{enumerate}[label=\textup{(\roman*)}]
    \item Fixing a common-difference set $S$ with $|S|\ge 2$ and searching for infinitely many supporting $n$.
    \item Any common-parameter Möbius / projective-rank-$2$ packet law.
    \item Any common one-parameter rational transfer law for $4$ rows.
    \item Treating large $\tau(n)$ alone as the right invariant.
\end{enumerate}
\end{observation}

\begin{observation}[what remains alive]
The live routes remaining in the notes are:
\begin{enumerate}[label=\textup{(\roman*)}]
    \item the anchored quartic / six square abscissae problem;
    \item backward square-translate construction of extendable row complexes;
    \item new width-$5$ source packets for the elliptic $3$-row law;
    \item the exact genus-$5$ $4$-row transfer surface;
    \item the moving non-split seed-centered two-parameter deformation.
\end{enumerate}
\end{observation}

\section{Collected open problems from the notes}

\begin{question}[anchored quartic problem]
Does there exist a quartic
\[
W^2=(T-\Delta_1)(T-\Delta_2)(T+\Delta_3)(T+\Delta_4)
\]
with at least six square integral abscissae $T$?
\end{question}

\begin{question}[spectral-cap problem]
Is there a universal bound
\[
|\Sigma(s_0;s_1,\dots,s_5)|\le 4
\]
for every anchored same-$\bmod 4$ sextuple, or does some anchored sextuple have spectrum size at least $5$?
\end{question}

\begin{question}[packet upgrade problem]
Does there exist a $3$-row packet with either
\begin{enumerate}[label=\textup{(\roman*)}]
    \item six positive common columns, or
    \item five positive common columns and a fourth row?
\end{enumerate}
The packet
\[
(c_2,c_3,c_4)=(756000,15971200,45130176)
\]
shows that width $5$ already occurs for $3$ rows, but it does not upgrade.
\end{question}

\begin{question}[genus-$5$ transfer problem]
Can one find
\[
0<c_2<c_3<c_4<c_5
\]
and five columns on the genus-$5$ transfer curve of Theorem 4.10, or prove a sharp obstruction for broad natural families?
\end{question}

\begin{question}[seed-centered deformation problem]
Does the moving two-parameter lift surface of Proposition 8.16 carry a rational point yielding a genuine lifted $4$-shift/$5$-square specialization?
\end{question}

\appendix

\section{Bibliographic record from the notes}

\begin{thebibliography}{99}

\bibitem{ER1997}
P.~Erd\H{o}s and M.~Rosenfeld,
\emph{The factor-difference set of integers},
Acta Arithmetica \textbf{79} (1997), 353--359.

\bibitem{JU1999}
M.~Jim\'enez-Urroz,
\emph{A note on a conjecture of Erd\H{o}s and Rosenfeld},
Journal of Number Theory \textbf{78} (1999), 140--143.

\bibitem{Bremner2019}
A.~Bremner,
\emph{On a problem of Erd\H{o}s related to common factor differences},
International Journal of Number Theory \textbf{15} (2019), 1059--1068.

\bibitem{CJ2000}
J.~Cilleruelo and M.~Jim\'enez-Urroz,
\emph{The hyperbola $xy=N$},
Journal de Th\'eorie des Nombres de Bordeaux \textbf{12} (2000), 87--100.

\bibitem{Chan2014}
S.~Chan,
\emph{Factors of a perfect square},
Acta Arithmetica \textbf{163} (2014), 171--180.

\bibitem{Chan2015}
S.~Chan,
\emph{Factors of almost squares and lattice points on circles},
arXiv:1406.2230.

\bibitem{Chan2025}
S.~Chan,
\emph{Numbers with three close factorizations and central lattice points on hyperbolas},
arXiv:2507.07094.

\bibitem{ChanEtAl2025}
S.~Chan, A.~Holmes, Y.~Liu, and J.~Villarreal,
\emph{Numbers with four close factorizations},
arXiv:2508.02818.

\bibitem{AlonEtAl}
N.~Alon, O.~Angel, I.~Benjamini, and E.~Lubetzky,
\emph{Sums and products along sparse graphs},
or related square-translate notes as cited in the session literature map.

\bibitem{DujellaElsholtz}
A.~Dujella and C.~Elsholtz,
\emph{Additive decompositions of the set of squares and related problems},
as cited in the session literature map.

\bibitem{EW2024}
C.~Elsholtz and A.~Wurzinger,
\emph{Unconditional upper bounds for sumsets contained in the squares},
as cited in the session literature map.

\end{thebibliography}

\end{document}
